\documentclass[12pt]{article}
\usepackage[a4paper, margin=1in]{geometry}
\usepackage{amsmath}
\usepackage{setspace}
\onehalfspacing

\title{Functions and Activities of the Department of Statistics: Kerala and India}
\author{Aleena Mary John}
\date{\today}

\begin{document}
\maketitle

\section*{Introduction}
Statistics departments are central to producing reliable data which governments, researchers and policy makers rely on for planning and decision making. In India, this is a joint responsibility of the national statistical machinery and the state level departments to ensure that data collection is locally relevant and nationally consistent. This report highlights the major functions and activities of the Department of Statistics in Kerala and the corresponding statistical set up at national level.

\section*{Department of Statistics, Kerala}
The Department of Economics and Statistics, Government of Kerala, is responsible for producing state-specific data that supports regional planning and welfare programs. Its major functions include:

\textbf{Agricultural Censuses:} The department conducts periodic agricultural censuses to track crop yields, patterns of land usage, and the costs involved in cultivation. This information is vital for assessing the productivity of the state's farming sector and for designing agricultural policy and subsidy programs suited to Kerala's unique cropping patterns.

\textbf{Economic Estimations:} A core activity of the department is the computation of the Gross State Domestic Product (GSDP), along with estimates of industrial growth and inflation indices. These figures allow the state government to monitor economic performance over time and to compare Kerala's growth trajectory with national trends.

\textbf{Vital Statistics:} The department also registers vital events such as births and deaths, and tracks broader demographic trends within the state. This data feeds into public health planning, resource allocation for hospitals and welfare schemes, and long-term population projections.

\textbf{Socioeconomic Surveys:} Through regular socioeconomic surveys, the department evaluates employment patterns and living standards across different districts and communities. These surveys help identify areas that require targeted intervention and are used to assess the effectiveness of ongoing welfare schemes.

\section*{Department of Statistics, India}
At the national level, statistical functions are coordinated primarily through the Ministry of Statistics and Programme Implementation (MoSPI) and its associated bodies, which perform the following key activities:

\textbf{Macroeconomic Data:} The national statistical system computes the country’s Gross Domestic Product (GDP), monitors inflation through indices such as the Consumer Price Index (CPI) and Wholesale Price Index (WPI), and gages industrial production through the Index of Industrial Production (IIP). Such indicators are a precondition for macroeconomic policy making such as fiscal and monetary policy decisions.

\textbf{National Surveys:} Large-scale national surveys are conducted to measure employment levels, household consumption patterns, and broader socio-economic metrics. These surveys are carried out by the National Sample Survey Office (NSSO), which provides a statistically representative picture of the country's population and its economic conditions.

\textbf{Demographic Tracking:} The Census of India is conducted every ten years by the national statistical machinery. The civil registration systems of vital statistics are also maintained. This provides the most complete demographic data set available for the country.

\textbf{Coordination:} A key function at the national level is aligning statistical standards, definitions, and methodologies across various ministries and state departments, which ensures the data collected across the country remains comparable and consistent.

\section*{Conclusion}
Together, the state and national statistical departments form an interconnected system: state departments like Kerala's capture region-specific detail, while the national framework ensures standardization and aggregation across the country. This division of responsibility allows India to maintain both granular, locally useful data and a coherent national statistical picture.

\end{document}